\documentclass{article}
\usepackage[UTF8]{ctex}
\usepackage{amsmath,amssymb,amsfonts}
\usepackage{enumitem}
\usepackage{geometry}
\geometry{a4paper, margin=1in}
% === 错题本排版优化（不参与题目标识） ===
\setlength{\parindent}{0pt}
\setlength{\parskip}{0.35em plus 0.10em minus 0.08em}
\setlength{\emergencystretch}{3em}
\setlist{itemsep=0.25em, topsep=0.40em, parsep=0pt}
\allowdisplaybreaks[2]
% === 错题本排版优化结束 ===
\begin{document}
\section*{中国海洋大学2024年数学分析真题}
\section*{124．中国海洋大学2024年数学分析试卷}

\subsection{题目1}
1 ．求极限 \(\lim _{x \rightarrow 0} \frac{1}{x}\left(\frac{1}{x}-\cos x\right)\) ．

\subsection{题目2}
2 ．利用 \(\frac{1}{n+1}<\ln \left(1+\frac{1}{n}\right)<\frac{1}{n}, n \in \mathbb{N}^{+}\)，计算 \(\lim _{n \rightarrow \infty} \frac{1+\frac{1}{2}+\frac{1}{3}+\cdots+\frac{1}{n}}{\ln n}\) ．

\subsection{题目3}
3．已知 \(\int_{0}^{+\infty} \frac{\sin x}{x}=\frac{\pi}{2}\) ，计算 \(\int_{0}^{+\infty}\left(\frac{\sin \alpha x}{x}\right)^{2} \mathrm{~d} x(\alpha>0)\) ．

\subsection{题目4}
1．判断正误：设 \(f(x)\) 是 \([a, b]\) 上的非负函数，且 \(\int_{a}^{b} f(x) \mathrm{d} x=0\) ．则在 \([a, b]\) 上一定存在非空开区间使得 \(f\) 恒为零。

\subsection{题目5}
2．\(f_{n}(x)=x^{n}-x^{3 n}, 0 \leq x \leq 1\) ，判断函数列 \(\left\{f_{n}(x)\right\}\) 在 \([0,1]\) 上的收敛性和一致收敛性．

\subsection{题目6}
3．判断正误：在 \([0,1]\) 上存在函数，使得在每个有理点处连续，但在每个无理点处不连续．

\subsection{题目7}
4．设 \(a_{n}>0\) 且级数 \(\sum_{n=1}^{\infty} a_{n}\) 发散，判断级数 \(\sum_{n=1}^{\infty} \frac{a_{n}}{1+n^{2} a_{n}}\) 的敛散性。

\subsection{题目8}
1．设 \(f(x)\) 为 \(U_{-}^{\circ}\left(x_{0}\right)\) 上的单调有界函数，证明 \(\lim _{x \rightarrow x_{0}^{-}} f(x)\) 存在．

\subsection{题目9}
2．叙述并证明拉格朗日中值定理．

四．设 \(f(x)\) 在 \([0,1]\) 上连续，且 \(f(0)=f(1)\) ．证明：对任何正整数 \(n\) ，必存在相应的 \(\xi \in(0,1)\) ，使得

\[
f\left(\xi+\frac{1}{n}\right)=f(\xi)
\]

五．设 \(f(x)\) 为 \([0,+\infty)\) 上正值单调减少函数，且 \(\int_{0}^{+\infty} f(x) \mathrm{d} x\) 收敛．证明：

\[
\lim _{h \rightarrow 0^{+}} h \sum_{n=1}^{\infty} f(n h)=\int_{0}^{+\infty} f(x) \mathrm{d} x
\]

六．计算第二类曲线积分

\[
I=\int_{L} \frac{(x+y) \mathrm{d} x-(x-y) \mathrm{d} y}{x^{2}+y^{2}} .
\]

其中 \(L\) 是从点 \(A(-1,0)\) 到点 \(B(1,0)\) 的一条在 \(x\) 轴上侧且不通过原点的光滑曲线．

七．求函数 \(f(x, y, z)=x^{4}+y^{4}+z^{4}\) 在条件 \(x y z=1\) 下的极值，并判断该极值是极大值还是极小值？为什么？

八．计算三重积分

\[
\iiint_{\Omega} z \mathrm{~d} x \mathrm{~d} y \mathrm{~d} z
\]

其中 \(\Omega\) 是球面 \(x^{2}+y^{2}+z^{2}=8\) 的上半部分与抛物面 \(z=\frac{1}{2}\left(x^{2}+y^{2}\right)\) 所围立体．
\end{document}
